\section{Introduction}
When language-model agents bid on behalf of competing firms, oversight must
address coordination that is invisible in the public record. Reading messages
may be impossible, and similar bids need not imply an agreement: agents can
respond to the same history or implement similar policies. Conversely, a firm
that learns a rival's private cost can produce suspicious dependence without
the rival's cooperation. A useful audit therefore asks both which pair is
dependent and what information could account for that dependence.

Pairwise net influence \citep{bonjour2022} compares a pair's statistical
influence with the strongest influence from a third agent. We test this screen
in a procurement game with a designated pair of LLM bidders. Private costs
provide an additional diagnostic: the partially observable statistic separates
action dependence from dependence of a bid on another firm's private input.
The distinction matters because the symmetric influence score can combine
either unilateral or reciprocal information flow. Our auditors read no
messages; this avoids reliance on message content, but does not constitute a
test of steganographic channels.

The contributions are:
\begin{enumerate}
\setlength{\itemsep}{0pt}
\item A pre-registered test of net influence on LLM auctions, with matched
honest controls, correlation and judge baselines, and separate measures of
discrimination, localisation, and abstention.
\item An audit comparison with and without private inputs, using an exact
decomposition into action dependence and directional leaks, plus a
partner-cost intervention on frozen contexts.
\item Targeted tests of honest look-alikes, unilateral leakage, cartel
placement, and model breadth, with explicit limits on what detection implies.
\end{enumerate}
The intended contribution is an assessment of the information an audit adds;
superiority to a simpler screen is an empirical question.
